\subsection{The disk gluing, its local links, and its cell counts}

We use the later surface $\mathcal X_2$ of cote 154/005 and /007.
Its faces are the ordinary relative positions, and its additional edges
are the sweeps across an old pseudoline. This convention is the one used
by the count on 154/120. The capped version $\mathcal X_0$ and the versions
with a supercritical face are distinguished below. All our arrangements
are finite, labelled, and of rank three. An old intersection may have
arbitrary multiplicity. The boundary of the disk obtained by cutting
along an added pseudoline is its two-sheeted normal-orientation cover.
Its generization packets and empty specialization fans give the
transition intervals of 154/035--038. We shorten exceptional intervals,
as on /038, so distinct closed transition intervals are disjoint.

We call a complementary interval an \emph{inertia interval}, or a corner
after it is contracted in the polygonal cell model. Ordinary gluing uses
the reflection selected on 154/039: the two extreme chord names are
exchanged. The same-chord endpoint correspondence rejected on that page
is a different gluing and is not used here. For a multiple vertex, the
whole fan is treated at once; the two extreme strands determine the two
endpoints, and the internal strands determine its weight.

\begin{resultenv}{T1}{Proved}{The lens lemma and the only repeated-vertex corner}
Let $D$ be ordinary. If two consecutive generization intervals on its
disk boundary correspond to the same old vertex, then $D$ passes through
exactly one old vertex. The intervening side either contains an ordinary
specialization interval or bounds an exceptional corridor beside an old
pseudoline. Consequently the complete transition cycle, including the
exceptional sweep intervals of $\mathcal X_2$, never has a corner bounded
by two generizations at the same vertex.
\end{resultenv}
\begin{proof}
Different vertices on $D$ use disjoint old-line sets. In the disk, the
chords in one such set are disjoint and form two opposite packets. Every
chord belonging to another such set crosses each of them, so that other
set has one packet on each of the two complementary boundary intervals.
Thus the two packets of one vertex cannot be consecutive if there is a
second vertex on $D$.

Suppose now that $D$ passes only through $v$. Let $C$ be the parallel
class of chords coming from the old lines through $v$. On either side,
let $P$ be the boundary arc between the two $C$-packets, and take the
outermost chord $g\in C$ facing $P$. The arc $P$ and $g$ bound a disk
$U$. Every other old chord has exactly one endpoint on $P$ and crosses
$g$ exactly once. Its part in $U$ is an arc from $P$ to $g$. No other
chord of $C$ lies in the interior of $U$.

Orient these arcs from $P$ to $g$. Their intersection graph has no
directed cycle. Here is the precise known result giving this assertion:
decone the rank-three arrangement $\Sigma\cup\{D\}$ at $g$ and use
$D$ as objective. The chosen lune lies in one affine sign region
$D>0$, after reversing its sign if necessary. On this sector $U$, the directed edges
of its oriented-matroid program are exactly the indicated orientations
(reverse the objective sign if necessary). Every strand here is
nonparallel to the objective: a strand through $v$ would belong to $C$.
The graph therefore has no undirected edges in $U$.
Sturmfels and Ziegler's definition of this graph on pp.~30--31,
Proposition~3.9 on p.~33, and the rank-three Euclidean result,
Corollary~4.5 on p.~40, imply the asserted absence of directed cycles
\cite{SturmfelsZiegler1993}. Restricting an acyclic directed graph to
the vertices and edges in $U$ preserves acyclicity.

If $U$ contains an old vertex, choose a source in this finite directed
graph. Along each incident strand, its incoming segment runs directly
from $P$ to that vertex: an earlier vertex would be a predecessor.
These incoming endpoints form a consecutive block on $P$. Indeed, a
strand beginning between the two extreme incoming endpoints and not
ending at the source must leave the intervening fan to reach $g$.
The Jordan separation theorem forces it to cross an incoming extreme
segment before the source, again giving a predecessor. The same
argument excludes a vertex inside the fan. Including every strand
through the source therefore gives an empty specialization fan with
base on $P$, hence an ordinary specialization interval between the
two generization packets.

If there is no such specialization interval, $U$ has no old vertex.
Its intervening strands are mutually disjoint, and the region between
$D$ and the old line underlying $g$ is exactly an empty corridor.
Conversely an empty corridor contains no old vertex that could be a
specialization apex on this side. The shortened exceptional transition
interval occupies its boundary side and separates the two generization
intervals. This proves the final assertion, including multiple vertices.
\end{proof}

\begin{resultenv}{T2}{Proved}{Ordinary corner circuits are four-cycles}
For the reflection gluing above, every vertex circuit using only
ordinary transition edges has four corners from four distinct relative
positions. It contains one gen--gen corner, one spec--spec corner, and
two mixed corners. This remains true when a position passes through
arbitrarily many old vertices or a semi-stable position has arbitrarily
many critical specializations.
\end{resultenv}
\begin{proof}
We first check a gen--gen corner. By T1 its two intervals belong to
distinct vertices $v,w$. Their old-line sets are disjoint. Choose
disjoint closed coordinate disks at $v,w$, containing no other old
vertex, in which the old pseudolines are radial strands and $D$ is a
diameter. The choice of one of the two lifted packets at each vertex
specifies one of its two resolutions. Denote the corresponding signs
by $\epsilon,\eta$. The four resulting extensions have the following
two localization coordinates; all their other coordinates agree:
\[
\begin{array}{c|c|c}
\text{position}&(\sigma(v),\sigma(w))&\text{corner type}\\\hline
D&(0,0)&G,G\\
D_v&(\epsilon,0)&S,G\\
D_{vw}&(\epsilon,\eta)&S,S\\
D_w&(0,\eta)&G,S.
\end{array}
\]
The resolutions commute because their supports are disjoint. Their
localizations differ, so the four labelled positions are distinct,
also modulo reorientation of the new pseudoline: a nonzero coordinate
outside the indicated zero set fixes that reorientation.

We verify that this square is the actual corner circuit, rather than
merely a commuting diagram of extensions. In a resolved fan the selected
generization packet becomes the base of the reverse specialization
fan. Its extreme strand names reverse order. The chosen reflection
exchanges those names, and thus identifies the same two \emph{spatial
sides} of the interval before and after resolution. On the selected
side each extreme incoming strand now encounters the old apex first.
It cannot acquire another specialization apex before it; every such
triangle would have to cross that first vertex. The portions of all
strands outside the resolving disk are unchanged. Hence the adjacent
inertia interval remains adjacent to the reverse transition on that
side. All possible changes at the other lifted packet lie outside
this inertia interval: before resolution that other packet was itself
a transition interval, and our two intervals were consecutive. This
argument applies to a fan of any multiplicity, using its two extremes.

An exceptional corridor cannot newly intrude into the tracked interval.
If one resolution produces a position near an old line through the
remaining zero vertex, the just resolved, different vertex lies in the
tracked lens and blocks emptiness on that side. Its empty corridor, if
present, is on the other side. If the old line does not contain the
remaining zero vertex, nearness to it is impossible. After both
resolutions, an exceptional wall in the tracked spec--spec gap would
have to contain both distinct apices $v,w$: these are the two fans
bounding that gap. No old line contains both, by the disjoint-support
argument. Thus no exceptional interval enters that gap either.
We may therefore perform the two resolutions in either order while
following the same corner. The local circuit is exactly the displayed
four-cycle. Its inverse endpoint maps return the starting germ to
itself, so no additional identification occurs after four steps.

It remains to prove that no other all-ordinary circuit has been missed.
The endpoint identifications form a finite permutation of corner germs,
so such a circuit is finite. Across an ordinary edge the integer
$\nu(D)$ changes by exactly one. Choose a position of maximal $\nu$
on the circuit. Both adjacent steps decrease $\nu$, so its corner is
gen--gen. T1 excludes a repeated-vertex gen--gen corner, and the preceding
local computation identifies its entire orbit with a four-cycle.
Thus every ordinary circuit has the stated form. Neither the argument
nor its conclusion bounds the total number of specializations at any
one position.
\end{proof}

\begin{resultenv}{T3}{Proved}{Exceptional completion, embeddings, and the surface counts}
The prescribed $\mathcal X_2$ is a well-defined compact connected
nonorientable surface, up to labelled cellular homeomorphism. Every
vertex has valence four, and every closed face, edge, and vertex embeds.
With $b=\sum_i n_i$ and $F(t)=\sum_j c_jt^j$ as above, its cell counts are
\[
 s_2=F(1),\qquad s_1=2F'(1)+2b,\qquad s_0=F'(1)+b.
\]
Its ordinary vertices are the gen--gen inertia intervals; its exceptional
vertices are the $2b$ flags $(v,e)$ of the original arrangement.
\end{resultenv}
\begin{proof}
We make the exceptional construction explicit. Fix an old pseudoline
$D_i$, with cyclically ordered vertices and edges. In a sufficiently
small Moebius neighborhood of $D_i$, an ordinary pseudoline close to
$D_i$ crosses it once. Its crossing is either in an old edge $e$, or at
an old vertex $v$. In each case there are two choices of the local
normal direction. We denote these positions $U_e^+,U_e^-$ and
$S_v^+,S_v^-$. The superscripts are local normal labels; transporting
them once around $D_i$ exchanges them. There are exactly $2n_i$ stable
positions $U_e^\pm$ and $2n_i$ semi-stable positions $S_v^\pm$ near
$D_i$. To verify exhaustiveness, cut the neighborhood at its single
crossing with $D_i$. On the resulting strip, a close pseudoline is a
section with fixed side until that crossing and the opposite side
after it. Contracting each old-vertex fan leaves no further choice.
Conversely these sections realize all the indicated positions.

If a position admits an empty corridor with $D_i$, it is one of these
positions. Slide its boundary through the corridor; there is no old
vertex in the interior, so no localization changes before reaching
a sufficiently small neighborhood. This is an explicit isotopy
relative to the old arrangement's combinatorial cells, and gives
the required comparison with the near positions. It introduces no
additional relative-position class.

The exceptional edges pair $U_e^+$ with $U_e^-$, and $S_v^+$ with
$S_v^-$. There is one such sweep for every unoriented edge and vertex
of $D_i$, hence $2n_i$ edges. At every original flag $v\in\partial e$
the four incident positions, in cyclic order, are
\[
 U_e^+,\quad S_v^+,\quad S_v^-,\quad U_e^-.
\]
The first and third sides are ordinary specialization/generization
sides, and the other two are exceptional sweeps. The endpoint rule
preserves the original flag $(v,e)$. On a stable sweep it preserves
the old vertex at that end; on a semi-stable sweep it preserves the
old edge on that side. These are the two equivalent descriptions of
the same strip gluing. The four positions are distinct: the stable
ones have no zero, the semi-stable ones have exactly $v$ as zero,
and the two normal choices differ on the neighboring fan signs.
A vertex off $D_i$ exists because the arrangement has rank three;
its fixed nonzero sign prevents these differences from being removed
by reorienting the new element. Thus each exceptional circuit is a
four-cycle with four distinct face labels.

There is no missing corner between two exceptional intervals. At a
stable near position, an exceptional interval is bounded by the two
ordinary fans at the distinct endpoints of its old edge. Two different
old wall lines cannot occupy that same gap: both would have to contain
those two vertices. At a semi-stable near position, the exceptional
interval lies between the two packets of its single vertex, and the
outermost wall of the lens on that side is unique. Thus every corner
meeting an exceptional edge is one of the flagged circuits just listed.

We also check exhaustiveness of the boundary circuits before inserting
any exceptional transition intervals or sweeps. The following maximum
argument is on those ordinary inertia circuits. A circuit not capped by T2 has a maximum of
$\nu$; its gen--gen corner must therefore use the same vertex, and
T1 identifies a specific empty corridor with an old line $D_i$.
Resolving its vertex toward an incident old edge $e$ gives $U_e$.
The corridor remains empty, and the next reverse fan is at the other
endpoint of $e$, giving $S_w$. Following the boundary thus moves along
the vertices and edges of the same $D_i$, without changing its label.
Conversely every such old flag supplies this local step. After one
turn the normal sign is reversed, because a pseudoline in the
projective plane is one-sided; after two turns it returns. This proves
that the boundary circuits are exactly the normal-orientation circles
of the old lines, with $4n_i$ successive stable/semi-stable sectors.

One can check the twisting without any orientation inference from the
counts. Orient all near positions as generizations of $[i^+]$, using
a nonzero cocircuit off $D_i$ to make this unique. Lemma~4.1 of
Sturmfels--Ziegler identifies their cyclic incidence with
$\mathcal E(M/i)$. The sweep sends
$\kappa_+(\rho)$ to $\kappa_+(-\rho)$: it reverses just the signs at
the vertices of $D_i$ and preserves the signs off $D_i$.
Thus the actual pairing of the free boundary intervals is the
antipodal identification on this circle. The quotient of its collar
by $(z,0)\sim(-z,0)$ is a Moebius band.

In explicit coordinates the inserted band is
$([0,1]\times[-1,1])/((0,t)\sim(1,-t))$.
Its boundary is that normal-orientation double cover. Put fiber edges
in the dual cell picture over the vertices and edge-midpoints of the
old $D_i$. The intervening
rectangles are indexed by the old flags and have exactly the four
corner labels displayed above. Thus the prescribed exceptional
endpoint pairing is the actual Moebius-band gluing, rather than a
conclusion drawn from an Euler-characteristic calculation.
Filling each such circuit by a disk gives the capped surface
$\mathcal X_0$. Replacing that disk by a Moebius band, with the flag
subdivision just described, gives exactly $\mathcal X_2$; this is the
topological operation stated on 154/007. In particular every boundary
component is filled, and the inserted Moebius band is an embedded
subsurface. This proves nonorientability directly, without relying on
the manuscript's monodromy assertion on 154/064.

For completeness, we check the topology and the embedding issue.
First glue the original disks only along their pairwise disjoint
closed transition intervals. Each interior point of such an interval
has exactly two representatives, and no interval endpoint belongs
to a second transition interval in the same disk. Thus no zigzag
of these identifications joins distinct points of one original disk.
The finite no-zigzag criterion SA10--A4.3, rechecked here, gives
embeddings. In the polygonal cell model, equivalence classes of corner
germs are exactly the four-cycles of T2 and the preceding exceptional
calculation. Each has four distinct disk labels. Hence no two corners
of a single normalized face are identified, no edge has its endpoints
identified, and no edge interior is identified with any other part
of the same face. The same no-zigzag criterion proves that all closed
cell characteristic maps embed. One may equivalently retain the
original disks and distribute the four sectors of each filled inertia
circuit among its incident faces; this gives the same closed-cell
model without collapsing an interval in the original disk.

The quotient relation on the finite union of closed polygonal disks is
closed: it is the finite union of paired edge graphs, diagonal interior
relations, and finite corner classes. The quotient is compact Hausdorff.
A face-interior point has a disk neighborhood, an edge-interior point
has two half-disks, and a vertex has the cone on a genuine four-cycle,
hence a disk neighborhood. This proves the surface assertion by actual
local links, as required by the distinction in SA10--C3.04.

The surface is connected. The cubical part of $T(M)$ obtained by
removing the interiors of the parallel polygonal cells is connected:
the full complex is connected by E1, and attaching a two-cell along
its connected boundary cannot join two components. Its face-incidence
graph is exactly the graph of ordinary positions and ordinary
transitions. That graph is therefore connected, and adding the
exceptional sweeps preserves connectedness.

There is one face per ordinary position, giving $s_2=F(1)$. A position
with $j$ zero vertices has exactly $2j$ immediate generizations,
each counted once at its more special endpoint. Thus there are
$2F'(1)$ ordinary edges. The exceptional calculation adds $2b$ edges.
Every vertex is four-valent, so $2s_1=4s_0$, giving the displayed
counts. Each ordinary four-cycle has exactly one gen--gen corner by
T2, and each exceptional vertex has exactly one original flag, proving
the final indexing assertion.

Finally, the construction depends only on the labelled chord diagrams
and the specified endpoint permutations. Changes of interval lengths
or parametrizations are absorbed by boundary homeomorphisms, which
extend across each disk by coning. Choice of a starting point or
orientation of a cyclic word merely reparametrizes the same polygon.
If the labelled internal chord graph is retained, match its vertices
and edges first, and extend across each complementary disk cell;
this gives the same extension while preserving that graph. These
extensions give a labelled cellular homeomorphism of the quotients.
This is the precise ambiguity in ``well-defined.''
\end{proof}

\paragraph{Scope of the edge cases.}
A rank-three arrangement of distinct pseudolines has $n\geq3$ and at
least two vertices on every old line. Thus neither an empty arrangement,
a one-line arrangement, nor a pencil satisfies the hypotheses. Multiple
intersections are included throughout T1--T3. All cells here are connected
disks or intervals; disconnected strata, infinite arrangements, and
non-locally-finite gluing families are not asserted to satisfy these
finite statements. No faithfulness or divisibility axiom from the
separate tame-space questions enters the construction.
